\section{Implementierung und Training}

\subsection{Entwicklungsumgebung}
Dieses Kapitel dokumentiert die verwendete Hardware, Software und Frameworks (Python, PyTorch) zur Sicherstellung der wissenschaftlichen Reproduzierbarkeit.

\subsection{Implementierung der Data Pipeline}

\subsubsection{Dataset-Klassen}
Dieser Abschnitt beschreibt die Programmierung der PyTorch-Schnittstelle zum Einlesen und Vorhalten der bereinigten Betriebsdaten.

\subsubsection{DataLoader und Batch-Generierung}
Es werden Laderoutinen zur iterativen Übergabe der Zeitreihenpakete (Batches) an das Modell umgesetzt.

\subsection{Trainingsschleife und Hyperparameter-Tuning}
Die programmierte Trainingsschleife wird dokumentiert, einschließlich Festlegung und Begründung von Learning Rates, Batch-Sizes und Early-Stopping-Kriterien.
